\chapter{انواع داده}
\label{ch:types}

\begin{goals}
\begin{itemize}
  \item چهار نوع داده‌ی پایه: عدد صحیح، عدد اعشاری، رشته و منطقی
  \item سلام چگونه نوع هر مقدار را تشخیص می‌دهد
  \item تقسیم عددهای صحیح و تفاوت آن با تقسیم اعشاری
  \item تبدیل نوع با \fcode{بعنوان}
  \item چسباندن متن و عدد به هم
\end{itemize}
\end{goals}

\section{هر مقدار یک نوع دارد}

در فصل پیش دیدیم که در جعبه‌ی یک متغیر می‌توان عدد یا متن گذاشت. اما رایانه
عدد و متن را یکسان نگه نمی‌دارد و یکسان با آن‌ها رفتار نمی‌کند. برای عددها
می‌توان جمع و ضرب حساب کرد، برای متن نه؛ متن را می‌توان حرف‌به‌حرف خواند،
عدد را نه. به همین دلیل هر مقدار در سلام یک \textbf{نوع} دارد و سلام
همیشه می‌داند در هر جعبه چه نوع چیزی است.

در این کتاب با چهار نوع پایه سروکار داریم:

\begin{center}
\begin{tabular}{lll}
\toprule
نوع & نام در سلام & نمونه \\
\midrule
عدد صحیح & \fcode{صحیح} & \code{17} ، \code{-4} ، \code{0} \\
عدد اعشاری & \fcode{اعشار۶۴} & \code{2.5} ، \code{-0.75} ، \code{3.0} \\
رشته (متن) & \fcode{رشته} & \code{"سلام"} ، \code{"17"} ، \code{""} \\
منطقی & \fcode{منطقی} & \fcode{درست} ، \fcode{نادرست} \\
\bottomrule
\end{tabular}
\end{center}

سلام از روی شکل مقداری که می‌نویسید، نوعش را خودش تشخیص می‌دهد. لازم
نیست نوع را بنویسید:

\program{چهار نوع داده}
\begin{salam}
تابع اصلی:
    تعداد := 12          // عدد صحیح
    قیمت := 2.5          // عدد اعشاری
    نام := "کتاب"        // رشته
    موجود := درست        // منطقی
    چاپ تعداد, قیمت, نام, موجود
تمام
\end{salam}

\begin{outputfa}
12 2.5 کتاب true
\end{outputfa}

قاعده‌ی تشخیص ساده است: عددی بدون نقطه‌ی اعشار، صحیح است؛ عددی با نقطه،
اعشاری است؛ هر چه میان نقل قول باشد، رشته است؛ و \fcode{درست} و
\fcode{نادرست} منطقی‌اند. به آخرین مقدار خروجی دقت کنید: سلام مقدارهای
منطقی را به شکل \code{true} و \code{false} چاپ می‌کند.

\section{عدد صحیح}

عدد صحیح همان عددهای بی‌اعشار است: تعداد دانش‌آموزان، شماره‌ی صفحه، سال.
می‌تواند منفی هم باشد. جمع، تفریق و ضرب دو عدد صحیح، همیشه یک عدد صحیح
می‌دهد. اما تقسیم داستان دیگری دارد:

\program{تقسیم عددهای صحیح}
\begin{salam}
تابع اصلی:
    چاپ 7 / 2
    چاپ 9 / 3
    چاپ 1 / 4
تمام
\end{salam}

\begin{outputfa}
3
3
0
\end{outputfa}

هفت تقسیم بر دو، سه و نیم است، اما سلام \code{3} چاپ کرد. چرا؟ چون هر دو
عدد صحیح بودند و \textbf{تقسیم دو عدد صحیح در سلام همیشه یک عدد صحیح
می‌دهد}: بخش اعشاری دور ریخته می‌شود. به همین دلیل یک تقسیم بر چهار، صفر
شد. این رفتار در آغاز عجیب به نظر می‌رسد، اما بسیار پرکاربرد است. اگر ۷
سیب را میان ۲ نفر تقسیم کنید، هر کس ۳ سیب کامل می‌گیرد و ۱ سیب باقی
می‌ماند. سلام برای همین «باقی‌مانده» هم علامتی دارد: \code{\%}.

\program{خارج قسمت و باقی‌مانده}
\begin{salam}
تابع اصلی:
    سیب := 7
    نفر := 2
    چاپ "سهم هر نفر:", سیب / نفر
    چاپ "باقی‌مانده:", سیب % نفر
تمام
\end{salam}

\begin{outputfa}
سهم هر نفر: 3
باقی‌مانده: 1
\end{outputfa}

عملگر باقی‌مانده در فصل بعد بارها به کارمان می‌آید؛ مثلاً برای فهمیدن این
که عددی زوج است یا فرد.

\begin{deep}[عدد صحیح تا کجا بزرگ می‌شود؟]
هر عدد صحیح در سلام در جعبه‌ای با اندازه‌ی مشخص نگه داشته می‌شود و به همین
دلیل حدی دارد: از حدود منفی دو میلیارد تا مثبت دو میلیارد. برای همه‌ی
برنامه‌های این کتاب این حد بیش از اندازه کافی است، اما اگر روزی خواستید
جمعیت جهان یا بودجه‌ی کشور را به ریال حساب کنید، به آن برمی‌خورید. سلام
برای این کارها نوع‌های بزرگ‌تر هم دارد، مثل \fcode{صحیح۶۴}، اما در این کتاب
به آن‌ها نیازی نخواهیم داشت.
\end{deep}

\section{عدد اعشاری}

وقتی به عددهایی با بخش اعشاری نیاز داریم (قیمت به دلار، وزن، معدل، دما)
از عدد اعشاری استفاده می‌کنیم. کافی است در نوشتن عدد نقطه بگذاریم. حتی
\code{3.0} هم یک عدد اعشاری است، هرچند مقدارش سه باشد:

\program{محاسبه با عددهای اعشاری}
\begin{salam}
تابع اصلی:
    چاپ 7.0 / 2.0
    چاپ 2.5 * 4.0
    چاپ 10.0 / 4.0
    چاپ 1.0 / 3.0
تمام
\end{salam}

\begin{outputfa}
3.5
10
2.5
0.3333333333333333
\end{outputfa}

دو نکته در این خروجی هست. نخست، \code{2.5 * 4.0} برابر \code{10.0} است، اما
سلام آن را \code{10} چاپ کرد: سلام هنگام چاپ عددهای اعشاری، صفرهای بی‌فایده‌ی
پایان عدد را نمی‌نویسد. دوم، یک سوم را با شانزده رقم اعشار چاپ کرد، چون
یک سوم اعشار بی‌پایان دارد و رایانه فقط تا حدی از آن را نگه می‌دارد.

\begin{warn}
عددهای اعشاری در رایانه \emph{تقریبی} نگه داشته می‌شوند. برنامه‌ی زیر را
اجرا کنید:
\begin{salam}
تابع اصلی:
    چاپ 0.1 + 0.2
تمام
\end{salam}
\begin{outputfa}
0.30000000000000004
\end{outputfa}
این اشکال سلام نیست؛ همه‌ی رایانه‌های جهان همین را می‌گویند، چون عدد ۰٫۱
در مبنای دو (که رایانه با آن کار می‌کند) اعشار بی‌پایان دارد، درست مثل یک
سوم در مبنای ده. برای کارهای روزمره این خطای بسیار کوچک هیچ اهمیتی
ندارد، اما به خاطر داشته باشید که هرگز نباید دو عدد اعشاری را برای
«دقیقاً برابر بودن» مقایسه کنید.
\end{warn}

\section{وقتی صحیح و اعشاری قاطی می‌شوند}

سلام اجازه نمی‌دهد یک عدد صحیح و یک عدد اعشاری را مستقیم با هم تقسیم یا
جمع کنید؛ باید اول هر دو را به یک نوع درآورید. برای این کار واژه‌ی
\kwd{بعنوان} هست که مقدار یک متغیر را به نوع دیگری تبدیل می‌کند:

\program{تبدیل صحیح به اعشاری}
\begin{salam}
تابع اصلی:
    مجموع نمره := 17
    تعداد درس := 4
    میانگین := (مجموع نمره بعنوان اعشار۶۴) / (تعداد درس بعنوان اعشار۶۴)
    چاپ "میانگین:", میانگین
تمام
\end{salam}

\begin{outputfa}
میانگین: 4.25
\end{outputfa}

اگر تبدیل را انجام نمی‌دادیم، \fcode{مجموع نمره / تعداد درس} تقسیم دو عدد
صحیح بود و \code{4} می‌داد. پرانتزها لازم‌اند تا روشن باشد که تبدیل روی
کدام مقدار انجام می‌شود.

تبدیل در جهت عکس هم ممکن است. وقتی عددی اعشاری را \fcode{بعنوان صحیح}
تبدیل می‌کنید، بخش اعشاری آن دور ریخته می‌شود (گرد نمی‌شود، بریده می‌شود):

\program{تبدیل اعشاری به صحیح}
\begin{salam}
تابع اصلی:
    قیمت دقیق := 3.75
    قیمت رند := قیمت دقیق بعنوان صحیح
    چاپ قیمت رند
    منفی := -3.75
    چاپ منفی بعنوان صحیح
تمام
\end{salam}

\begin{outputfa}
3
-3
\end{outputfa}

\begin{note}
\fcode{بعنوان} روی یک \emph{متغیر} کار می‌کند. اگر بخواهید عددی را که همان
جا نوشته‌اید تبدیل کنید، بهتر است اول آن را در متغیری بگذارید و سپس تبدیل
کنید؛ همان کاری که در برنامه‌ی بالا با \fcode{منفی} کردیم.
\end{note}

\section{رشته}

رشته یعنی متن: دنباله‌ای از حرف‌ها، رقم‌ها و علامت‌ها که میان نقل قول
نوشته می‌شود. رشته می‌تواند خالی باشد (\code{""}) و می‌تواند رقم داشته باشد،
اما رقمِ درون رشته عدد نیست:

\program{رشته یا عدد؟}
\begin{salam}
تابع اصلی:
    چاپ 12 + 3
    چاپ "12" + 3
    چاپ "12" + "3"
تمام
\end{salam}

\begin{outputfa}
15
123
123
\end{outputfa}

خط دوم و سوم را ببینید. علامت \code{+} برای رشته‌ها معنی «چسباندن» دارد،
نه جمع کردن. \fcode{"12" + 3} یعنی رشته‌ی «۱۲» را به عدد ۳ بچسبان، که
می‌شود رشته‌ی «۱۲۳». اگر یکی از دو طرف \code{+} رشته باشد، سلام طرف دیگر را
هم به رشته تبدیل می‌کند و دو رشته را به هم می‌چسباند. این ویژگی برای ساختن
پیام‌ها بسیار مفید است:

\program{ساختن پیام}
\begin{salam}
تابع اصلی:
    نام := "سارا"
    سن := 17
    پیام := نام + " " + سن + " سال دارد."
    چاپ پیام
تمام
\end{salam}

\begin{outputfa}
سارا 17 سال دارد.
\end{outputfa}

به فاصله‌ها دقت کنید: چسباندن، خودش فاصله‌ای نمی‌گذارد، پس هر جا فاصله
می‌خواستیم، رشته‌ای با یک فاصله (\code{" "}) اضافه کردیم. مقایسه کنید با
\fcode{چاپ نام, سن} که میان چیزهایش خودش فاصله می‌گذارد. هر دو راه درست‌اند؛
با \code{+} کنترل بیشتری بر شکل پیام دارید و با ویرگول راحت‌ترید.

اگر رشته‌ای رقم دارد و می‌خواهید با آن حساب کنید، باید آن را به عدد تبدیل
کنید. برای این کار رشته‌ها یک قابلیت به نام \fcode{به صحیح} دارند که با
نقطه به آن‌ها می‌چسبد:

\program{تبدیل رشته به عدد}
\begin{salam}
تابع اصلی:
    متن := "12"
    عدد := متن.به صحیح()
    چاپ عدد + 3
تمام
\end{salam}

\begin{outputfa}
15
\end{outputfa}

به همین شکل \fcode{به اعشار} یک رشته را به عدد اعشاری تبدیل می‌کند. در فصل
\ref{ch:strings} با قابلیت‌های دیگر رشته‌ها آشنا می‌شویم.

\section{منطقی}

نوع منطقی فقط دو مقدار دارد: \fcode{درست} و \fcode{نادرست}. شاید کوچک به
نظر برسد، اما همه‌ی تصمیم‌گیری‌های برنامه (فصل \ref{ch:conditions}) بر
پایه‌ی همین دو مقدار است. بیشتر وقت‌ها مقدار منطقی را مستقیم نمی‌نویسیم،
بلکه از یک مقایسه به دست می‌آید:

\program{مقدارهای منطقی}
\begin{salam}
تابع اصلی:
    سن := 20
    بزرگسال است := سن >= 18
    چاپ بزرگسال است
    چاپ سن == 18
    چاپ "سلام" == "سلام"
تمام
\end{salam}

\begin{outputfa}
true
false
true
\end{outputfa}

علامت \code{>=} یعنی «بزرگ‌تر یا مساوی» و علامت \code{==} (دو مساوی پشت
هم) یعنی «آیا برابر است؟». دقت کنید که \code{==} با \code{=} فرق دارد: اولی
می‌پرسد، دومی مقدار می‌دهد. در فصل بعد همه‌ی علامت‌های مقایسه را می‌بینیم.

\section{نوع یک متغیر عوض نمی‌شود}

وقتی متغیری با یک نوع ساخته می‌شود، تا آخر همان نوع را دارد. نمی‌توان در
جعبه‌ای که برای عدد ساخته شده، متن گذاشت:

\begin{salam}
تابع اصلی:
    متغیر مقدار := 10
    مقدار = "ده"
    چاپ مقدار
تمام
\end{salam}

\begin{outputfa}
خطا: نمی‌توان 'str' را به 'i32' نسبت داد
 --> 3:5
\end{outputfa}

در پیام خطا، \code{str} نام انگلیسی نوع رشته و \code{i32} نام انگلیسی نوع
صحیح است؛ گاهی این نام‌ها را در پیام‌های خطا می‌بینید. این سخت‌گیری هم مثل
بقیه‌ی سخت‌گیری‌های سلام به سود شماست: وقتی متغیری به نام \fcode{تعداد}
می‌بینید، مطمئن هستید که همیشه یک عدد در آن است.

\section{یک برنامه‌ی کامل‌تر}

محاسبه‌ی معدل سه درس با واحدهای متفاوت، که همه‌ی نوع‌های این فصل را کنار
هم می‌آورد:

\program{معدل وزنی}
\begin{salam}
تابع اصلی:
    // نمره‌ها اعشاری‌اند، واحدها صحیح
    نمره ریاضی := 17.5
    واحد ریاضی := 3
    نمره فیزیک := 15.0
    واحد فیزیک := 2
    نمره ادبیات := 18.25
    واحد ادبیات := 2

    مجموع واحد := واحد ریاضی + واحد فیزیک + واحد ادبیات
    مجموع وزنی := نمره ریاضی * (واحد ریاضی بعنوان اعشار۶۴) +
                  نمره فیزیک * (واحد فیزیک بعنوان اعشار۶۴) +
                  نمره ادبیات * (واحد ادبیات بعنوان اعشار۶۴)
    معدل := مجموع وزنی / (مجموع واحد بعنوان اعشار۶۴)

    چاپ "مجموع واحدها:", مجموع واحد
    چاپ "معدل:", معدل
    چاپ "قبول؟", معدل >= 12.0
تمام
\end{salam}

\begin{outputfa}
مجموع واحدها: 7
معدل: 17
قبول؟ true
\end{outputfa}

به دو چیز دقت کنید. نخست این که یک محاسبه‌ی طولانی را در چند خط شکسته‌ایم؛
وقتی خطی به عملگری مثل \code{+} ختم می‌شود، سلام می‌فهمد که عبارت هنوز
تمام نشده و آن را در خط بعد دنبال می‌کند. دوم این که معدل واقعاً ۱۷٫۰ است (مجموع وزنی ۱۱۹ تقسیم بر ۷) و
سلام مثل همیشه آن را \code{17} چاپ کرده است.

\begin{summary}
\begin{itemize}
  \item هر مقدار یک نوع دارد: \fcode{صحیح} (عدد بی‌اعشار)، \fcode{اعشار۶۴}
        (عدد با اعشار)، \fcode{رشته} (متن) و \fcode{منطقی} (درست یا
        نادرست). سلام نوع را از شکل مقدار تشخیص می‌دهد.
  \item تقسیم دو عدد صحیح، بخش اعشاری را دور می‌ریزد. \code{\%}
        باقی‌مانده‌ی تقسیم را می‌دهد.
  \item عددهای اعشاری تقریبی‌اند؛ آن‌ها را برای برابری دقیق مقایسه نکنید.
  \item \fcode{مقدار بعنوان نوع} یک متغیر را به نوع دیگری تبدیل می‌کند.
        تبدیل اعشاری به صحیح، اعشار را می‌برد.
  \item \code{+} میان رشته و هر چیز دیگر، به معنی چسباندن است.
        \fcode{متن.به صحیح()} رشته را به عدد تبدیل می‌کند.
  \item نوع یک متغیر بعد از ساخته شدن عوض نمی‌شود.
\end{itemize}
\end{summary}

\section*{تمرین‌ها}

\exercise بدون اجرا کردن بگویید هر یک از این دستورها چه چاپ می‌کند و سپس
بررسی کنید: \fcode{چاپ 10 / 3}، \fcode{چاپ 10.0 / 3.0}، \fcode{چاپ 10 \% 3}،
\fcode{چاپ "10" + 3}، \fcode{چاپ 10 + 3}.

\exercise دمای هوا به سلسیوس را در متغیری بگذارید و آن را به فارنهایت تبدیل
و چاپ کنید. فرمول: فارنهایت برابر است با سلسیوس ضرب در ۹ تقسیم بر ۵،
به علاوه‌ی ۳۲. برنامه را یک بار با عدد اعشاری و یک بار با عدد صحیح
بنویسید و تفاوت خروجی‌ها را توضیح دهید.

\exercise تعداد کل دقیقه‌های یک فیلم (مثلاً ۱۳۵) را در متغیری بگذارید و
با تقسیم و باقی‌مانده، آن را به شکل «۲ ساعت و ۱۵ دقیقه» چاپ کنید.

\exercise سه متغیر با مقدارهای \code{"3"}، \code{3} و \code{3.0} بسازید.
آن‌ها را چاپ کنید، سپس هر یک را با عدد ۲ جمع کنید و نتیجه را چاپ کنید.
پیش از اجرا، خروجی را حدس بزنید.

\exercise قیمت یک کالا به دلار (اعشاری) و نرخ دلار به تومان (صحیح) را در
دو متغیر بگذارید و قیمت کالا به تومان را به شکل یک عدد صحیح چاپ کنید.
